This section preserves the first-release measurements and their original
source snapshots. The later implementation and its separately saved runs
are described in Section~\ref{sec:completion}; the figures below are not
reruns of the updated solver.

The original pruned-grid diagnostic implemented two-pass tree messages,
all coordinate min-marginals, geometric grids, interval filtering, and
hull recentering. Its saved exact-rational record covers nine fixtures,
70 stages, 26,589 bag-table states, and 526 pruned intervals. Separate
finite enumeration checked min-marginals on 21,034 assignments. These
figures describe that diagnostic, not the continuation's benchmark suite.
It did not implement the unknown-growth caps, rational reconstruction,
or an independent stored-history verifier. Its successes establish
finite regression evidence, not competitive general solver performance.

The continuation distinguishes three empirical questions: whether the
produced certificate independently checks, whether filtering reduces
state counts against comparable unfiltered methods, and whether its
total time and memory are useful against other solvers. An ordinary
floating-point solver result is useful as a timing baseline but is not
automatically an exact global certificate. Controlled synthetic fixtures
and public-library instances must be reported separately, with explicit
limits and incomplete runs preserved.

\subsection{The initial sparse solver and independent replay}
The initial standard-library Python reference implementation handled supplied
general tree decompositions, branching trees, empty separators, continuous
and native-integer domains, fixed coordinates, and disconnected factor
graphs. It implements corrected geometric grids, all coordinate
min-marginals, domain filtering, and the capped unknown-growth
approximation schedule. Coordinate-specific diagonal corrections sharpen
the common-$L$ bound. Endpoint reduction, bounded exact coordinate
descent, and an optional rational PSD/KKT presolve supplement the grid
algorithm. A saved result contains original-model bounds and a complete
history, including when a resource limit interrupts refinement.

The separate verifier does not call grid generation, dynamic programming,
filtering, coordinate descent, or the convex presolve. It shares rational
model parsing and objective evaluation, independently reconstructs factor
tables and Bellman equalities, checks all min-marginals and exclusions,
and re-evaluates feasible incumbents. The certificate proves the embedded
rational model; a caller must also match that model to the intended input.
It is an independently implemented checker, not a formally verified one.
Replay can cost as much as solving because it recomputes finite-table
arithmetic. This first release returned certified approximations and certain
exact presolve or finite-grid answers. General continuous rational
reconstruction was a theorem at that stage; the completion described below
adds its implementation.

\subsection{Bounded comparison on original objectives}
The main suite has ten small rational instances and two unmodified binary
QPLIB instances. The small set includes random signed paths, bands and a
branching tree, mixed and integer paths, flat and tied optima, and two
positive-curvature scales. Random cases do not have planted optimizers;
deliberate structural cases are labelled diagnostics. An independent exact
integer-slice and active-face enumeration gives the small reference values.

All twelve cases use absolute tolerance $1/50$; a random path and both
curvature cases also use $1/1000$. The resulting 15 configurations are
run with four methods. Each has a two-second solver limit, with an
external 14-second worker limit including setup, serialization and replay.
Exact methods additionally have a 20,000 aggregate bag-state cap and
24-stage cap. SCIP is single-threaded with a 256 MB internal limit.
The grid methods share a greedy minimum-degree decomposition; SCIP does
not receive it. These timings come from one bounded run on a shared
machine, not a statistical performance study.

\begin{center}
\begin{tabular}{lrrl}
\toprule
Method & Small targets reached & Public targets reached & Incomplete runs\\
\midrule
Geometric, pruned & $13/13$ & $0/2$ & 2 table limits\\
Geometric, unpruned & $13/13$ & $0/2$ & 2 table limits\\
Uniform, unpruned & $11/13$ & $0/2$ & 4 table limits\\
SCIP (numerical bounds) & $13/13$ & $0/2$ & 2 time limits\\
\bottomrule
\end{tabular}
\end{center}

``Target reached'' means the requested objective tolerance, not exact
optimization. The SCIP small outcomes comprise six optimal statuses and
seven gap-limit statuses. Its dual bounds remain numerical; its repaired
box-feasible incumbents are separately evaluated exactly. All 45 rational
grid certificates, including incomplete runs, pass replay. All 39
small-instance grid enclosures contain their exact reference values.

For the 13 matched small configurations, the pruned geometric method
processes 6,167 cumulative bag-table states versus 27,386 without pruning.
The strongest reduction in this small set is the high-positive-curvature
case at tolerance $1/1000$: 771 versus 13,527 states. Tied-and-flat and
pure-integer cases have equal counts for the two geometric variants.
Thus pruning helps particular fixtures, while the suite does not show
universal pruning benefit. Median solve times on these 13 configurations
are approximately 0.0059 seconds with pruning and 0.0082 without it;
median replay times are 0.0084 and 0.0078 seconds, respectively. Median
whole-worker subprocess times are 0.174 and 0.193 seconds. Interpreter
startup, setup, exact arithmetic and checking materially exceed some solve
times. Peak worker memory, including replay, is about 25 MB on the small
grid runs; it is not isolated DP memory.

The public inputs are QPLIB\_3852 (231 binary variables, 440 quadratic
terms) and QPLIB\_5881 (120 variables, 2,123 terms). Their supplied
decomposition widths are 25 and 95, respectively, and are heuristic
upper bounds on treewidth. Each exact variant refuses its oversized
tables before allocation and retains a valid initial bound; SCIP reaches
its time limit. No public instance reaches the target. These outcomes
demonstrate a practical width barrier for the supplied representations;
they do not prove that better decompositions or formulations cannot work.
The benchmark validates QPLIB's triangular coefficient convention against
the bundled library solutions and preserves the original objectives and
domains, apart from sign reversal of a maximization objective.

\subsection{Dimension and certified affine reduction}
Fourteen additional runs test increasing path dimension, a width-two
case, and the affine-recourse reduction. Across both first-release suites,
all \textbf{56 rational certificates} pass independent replay; the four
affine enclosures also contain their analytic zero optimum. The main
and extension suites use 20.23 and 22.23 seconds of summed subprocess
wall time, respectively. No worker crashes or reaches its external wall
limit. The following timings keep solve and replay separate.

\begin{center}
\begin{tabular}{lrrrr}
\toprule
Pruned random path & Gap & Solve (s) & Replay (s) & Peak MiB\\
\midrule
$n=16$ & $0.01497$ & $0.0253$ & $0.0231$ & $24.3$\\
$n=64$ & $0.01361$ & $0.2020$ & $0.1967$ & $24.8$\\
$n=256$ & $0.05786$ & $2.0023$ & $1.2388$ & $46.5$\\
\bottomrule
\end{tabular}
\end{center}

The target is $1/50$. The first two paths finish; the 256-variable
pruned path hits its time limit. Its unpruned counterpart hits a table
limit at the same gap and uses 1.790 seconds for replay. SCIP reaches
a numerical gap of approximately $0.000612$ in 1.797 seconds on that
path. A pruned width-two, 16-variable case also reaches its target, with
gap $0.01953$, solve time $0.0318$ and replay time $0.0546$ seconds.
Three paths and one wider example do not constitute an empirical
complexity law, and width one does not ensure fast Python rational
arithmetic at every dimension.

The affine experiment has six original variables, four retained
variables, and an exactly checked identity
$F(u,v,y)=q(u,v)+M\sum_i(y_i-u_i)^2$ with feasible response $y=u$.
It is the companion's designed two-block example. The same reduced
rational core is solved once for all three stiffness values.

\begin{center}
\begin{tabular}{lrrrr}
\toprule
Representation & States & Solve (s) & Replay (s) & Gap\\
\midrule
Original, $M=1$ & $1{,}694$ & $0.0352$ & $0.0425$ & $0.000458$\\
Original, $M=100$ & $64{,}564$ & $2.0044$ & $2.0668$ & $0.399817$\\
Original, $M=1000$ & $76{,}986$ & $2.0033$ & $2.5630$ & $7.406978$\\
Reduced, all $M$ & $394$ & $0.0121$ & $0.0151$ & $0.000519$\\
\bottomrule
\end{tabular}
\end{center}

The target is $1/1000$. The stiff original cases time out; the reduced
case finishes and its rational incumbent and lower bound lift exactly
to each original instance. State counts are cumulative over completed
stages. This demonstrates a useful certified
preprocessing operation on the stated class. It does not show that
these analytically simple models are difficult for other solvers, and
the affine-selector recognizer itself is not benchmarked here.

\subsection{Initial extension implementations}
The first-release conditional-cut implementation included residual sign
recognition, an exact rational max-flow oracle, primal/flow certificates,
adaptive core cells, trace verification, and a rational-separation exact
mode. Its broader concave/convex submodular extension had small exact
diagnostics. The TU implementation was a bounded two-bag diagnostic;
affine recourse and optimal-set/boundary extensions also had exact
algebraic checks. General convex-QP factor evaluation, diagonal-class LP
recovery, and reusable constrained-tree search were implementation gaps
at that stage. Section~\ref{sec:completion} records which gaps the next
release closes and states the remaining limits.

The saved experiment files provide per-run times, peak process memory,
decomposition metadata, source hashes, rational reference values and
replay certificates. The report's coverage map points to them. The evidence
supports a reusable certifying reference solver and concrete extensions;
it does not establish competitive general-purpose MINLP performance.
